% !TEX program = xelatex
% Deck EXPÉRIMENTAL — précision de l'estimation de Jaccard.
% Compiler avec XeLaTeX (Metropolis + Fira).  Utilise pgfplots.
% Contenu fondé sur la revue de C. Marchet (kamimrcht.github.io/webpage/sketch.html)
% et les articles qui y sont cités (Broder 1997 ; Li & Konig 2010 ; Flajolet+ 2007 ;
% Ondov+ 2016 / Mash ; Yu & Weber 2017 / HyperMinHash ; Baker & Langmead 2019).
\PassOptionsToPackage{table}{xcolor}
\documentclass[aspectratio=169,11pt]{beamer}
\usetheme[progressbar=frametitle, numbering=fraction, sectionpage=progressbar, block=fill]{metropolis}

% Police Fira par nom de fichier via kpathsea (portable).
\usepackage{fontspec}
\setsansfont{FiraSans}[
  Extension      = .otf,
  UprightFont    = *-Regular, ItalicFont = *-Italic,
  BoldFont       = *-SemiBold, BoldItalicFont = *-SemiBoldItalic,
]
\setmonofont{FiraMono}[Extension=.otf, UprightFont=*-Regular, BoldFont=*-Bold, Scale=MatchLowercase]

\usepackage[french]{babel}
\usepackage{csquotes}
\usepackage{amsmath,amssymb}
\usepackage{tikz}
\usetikzlibrary{positioning,arrows.meta,calc,decorations.pathreplacing,babel}
\usepackage{pgfplots}
\pgfplotsset{compat=1.18}

\definecolor{cBlue}{HTML}{1F6FB2}
\definecolor{cOrange}{HTML}{EB811B}
\definecolor{cGreen}{HTML}{2E8B57}
\definecolor{cRed}{HTML}{C0392B}
\definecolor{cGray}{HTML}{9AA0A6}

\tikzset{arr/.style={-{Stealth[length=2.3mm]}, line width=0.5pt}}

\title{Précision de l'estimation de Jaccard}
\subtitle{annexe expérimentale : variance, loi en $1/\sqrt{s}$, distance de Mash}
\author{TP \textsc{Jc2bimmm}}
\date{}
\institute{d'après la revue de C.~Marchet et les articles cités}

\begin{document}

\maketitle

%% ----------------------------------------------------------------
\begin{frame}{La question : quelle confiance dans $\hat J$ ?}
  \footnotesize
  On remplace des millions de k-mers par \textbf{$s$ valeurs}, puis on estime
  $\hat J = x/s$ ($x$ = nombre de positions \emph{concordantes}).
  \medskip

  Avec le \textbf{partition sketch}, les $s$ cases sont \textbf{indépendantes} :
  chaque case concorde avec probabilité $J$ (le vrai Jaccard). Autant de tirages
  \enquote{pile ou face}.

  \begin{center}
  \begin{tikzpicture}
    \foreach \i/\c in {0/{\textcolor{cGreen}{\checkmark}}, 1/{\textcolor{cGreen}{\checkmark}},
                       2/{\textcolor{cRed}{$\times$}}, 3/{\textcolor{cGreen}{\checkmark}},
                       4/{\textcolor{cRed}{$\times$}}, 5/{\textcolor{cGreen}{\checkmark}},
                       6/{\textcolor{cGreen}{\checkmark}}, 7/{\textcolor{cRed}{$\times$}}}
      \node[draw,minimum size=6.5mm] at (\i*0.8,0){\c};
    \draw[decorate,decoration={brace,mirror,amplitude=4pt}] (-0.35,-0.45)--(6.0,-0.45)
      node[midway,below,font=\scriptsize]{$s$ cases};
    \node[anchor=west,font=\scriptsize] at (6.4,0){$x=5 \Rightarrow \hat J = 5/8$};
  \end{tikzpicture}
  \end{center}
  \centering\footnotesize\textcolor{cGray}{De combien $\hat J$ s'écarte-t-il de $J$ ?}
\end{frame}

%% ----------------------------------------------------------------
\begin{frame}{Le modèle binomial}
  \footnotesize
  Le nombre de concordances $X$ suit une loi \textbf{binomiale} :
  \[ X \sim \mathcal{B}(s, J), \qquad \hat J = \frac{X}{s}. \]
  \begin{columns}[T]
    \begin{column}{0.54\textwidth}
      \textbf{Sans biais :} $\mathbb{E}[\hat J] = J$.
      \medskip

      \textbf{Variance :}
      \[ \operatorname{Var}(\hat J) = \frac{J(1-J)}{s},
         \quad \sigma = \sqrt{\frac{J(1-J)}{s}}. \]
    \end{column}
    \begin{column}{0.46\textwidth}
      \begin{alertblock}{Loi de l'erreur}
        \footnotesize
        \[ \varepsilon = O\!\left(\frac{1}{\sqrt{s}}\right) \]
        Indépendante de la taille du génome \textit{(Mash)}.
      \end{alertblock}
    \end{column}
  \end{columns}
  \medskip
  \centering
  \textcolor{cGreen}{\textbf{Diviser l'erreur par 2 $\Rightarrow$ multiplier $s$ par 4.}}
\end{frame}

%% ----------------------------------------------------------------
\begin{frame}{L'erreur décroît en $1/\sqrt{s}$}
  \centering
  \begin{tikzpicture}
    \begin{axis}[
      width=10.2cm, height=5.4cm,
      xmode=log, ymode=log, log basis x=2,
      xlabel={taille du sketch $s$}, ylabel={écart-type $\sigma(\hat J)$},
      xmin=16, xmax=8192, ymin=0.004, ymax=0.2,
      grid=both, grid style={gray!18},
      tick label style={font=\scriptsize}, label style={font=\footnotesize},
      legend cell align=left, legend style={font=\scriptsize, at={(0.97,0.97)}, anchor=north east},
    ]
      \addplot[cBlue,  very thick, samples=80, domain=16:8192] {sqrt(0.25/x)};
      \addlegendentry{$J=0{,}5$ \;(variance max)}
      \addplot[cOrange,very thick, samples=80, domain=16:8192] {sqrt(0.09/x)};
      \addlegendentry{$J=0{,}1$ ou $0{,}9$}
    \end{axis}
  \end{tikzpicture}

  \footnotesize\textcolor{cGray}{En échelle log--log, une \textbf{droite de pente $-\tfrac12$} :
  $\sigma \propto 1/\sqrt{s}$. La variance est maximale à $J=0{,}5$.}
\end{frame}

%% ----------------------------------------------------------------
\begin{frame}{Intervalle de confiance \& limites}
  \footnotesize
  \textbf{IC à 95\,\%} (approximation normale) :
  \[ \hat J \;\pm\; 1{,}96\,\sqrt{\frac{\hat J\,(1-\hat J)}{s}}. \]
  \begin{itemize}
    \item La variance \emph{absolue} est maximale à $J=0{,}5$ ; mais en \textbf{relatif},
          ce sont les \textbf{petits $J$} (génomes éloignés) qui souffrent : signal faible
          $\Rightarrow$ il faut un $s$ plus grand.
    \item \textcolor{cGray}{Rappel (revue) : le proxy Jaccard se dégrade quand les ensembles
          sont \emph{dissemblables} ou de \emph{tailles très différentes} $\Rightarrow$ on lui
          préfère alors la \emph{containment} (Jaccard de confinement).}
  \end{itemize}
\end{frame}

%% ----------------------------------------------------------------
\begin{frame}{De Jaccard à une distance : Mash}
  \footnotesize
  Mash convertit le Jaccard estimé $j$ en une distance proche de $1-\text{ANI}$ :
  \[ \boxed{\,D = -\frac{1}{k}\,\ln\!\frac{2j}{1+j}\,}
     \qquad \text{\scriptsize(Ondov \textit{et al.} 2016)} \]
  \begin{columns}[T]
    \begin{column}{0.56\textwidth}
      \textbf{Modèle de Poisson :} la probabilité qu'un k-mer \emph{échappe} à toute
      mutation (taux $d$) vaut $e^{-k d}$ — d'où le lien entre fraction de k-mers
      partagés et distance.
      \smallskip

      L'erreur sur $j$ (en $1/\sqrt{s}$) se \textbf{propage} à $D$.
    \end{column}
    \begin{column}{0.44\textwidth}
      \begin{exampleblock}{Exemple chiffré (Mash)}
        \footnotesize
        $s = 1000$, $D = 0{,}05$ :\\
        erreur $< 0{,}0068$ avec probabilité $0{,}99$.
      \end{exampleblock}
    \end{column}
  \end{columns}
\end{frame}

%% ----------------------------------------------------------------
\begin{frame}{Compresser : que coûte la précision ?}
  \footnotesize
  \begin{itemize}
    \item \textbf{b-bit MinHash} (Li \& König 2010) : ne garder que $b$ bits par haché
          $\Rightarrow$ collisions \emph{aléatoires} (proba $\approx 2^{-b}$) ; un estimateur
          \textbf{corrigé} retire ce biais, au prix d'un peu de variance — pour beaucoup
          moins de mémoire.
    \item \textbf{HyperMinHash} (Yu \& Weber 2017) : exposant \textsc{hll} (zéros de tête,
          $O(\log\log n)$ bits) $+$ mantisse $b$-bit $\Rightarrow \sim 16$ bits/case ;
          erreur relative $\varepsilon$ avec $\sim \varepsilon^{-2}$ cases.
    \item \textbf{HyperLogLog} (Flajolet \textit{et al.} 2007) : cardinalité à
          $\approx 1{,}04/\sqrt{m}$ ($m$ registres).
  \end{itemize}
  \begin{alertblock}{Compromis}
    \footnotesize
    À \textbf{mémoire égale}, plus de cases $\Rightarrow$ variance plus faible.
    Les collisions de la compression $\Rightarrow$ léger \textbf{biais} (surestimation),
    atténué par de plus gros sketches.
  \end{alertblock}
\end{frame}

%% ----------------------------------------------------------------
\begin{frame}{Références}
  \footnotesize
  \begin{itemize}
    \item \textbf{Broder} (1997). \textit{On the resemblance and containment of documents} — MinHash.
    \item \textbf{Li \& König} (2010). \textit{b-Bit Minwise Hashing} — WWW.
    \item \textbf{Flajolet, Fusy, Gandouet, Meunier} (2007). \textit{HyperLogLog} — AofA.
    \item \textbf{Ondov \textit{et al.}} (2016). \textit{Mash: fast genome and metagenome distance
          estimation using MinHash} — Genome Biology.
    \item \textbf{Yu \& Weber} (2017). \textit{HyperMinHash} — arXiv:1710.08436.
    \item \textbf{Baker \& Langmead} (2019). \textit{Dashing: fast and accurate genomic distances
          with HyperLogLog} — Genome Biology.
  \end{itemize}
  \vfill
  {\scriptsize\textcolor{cGray}{Synthèse : revue de C.~Marchet,
  \texttt{kamimrcht.github.io/webpage/sketch.html}.}}
\end{frame}

\begin{frame}[standout]
  La précision se résume à une loi :\\[2pt]
  \alert{$\sigma(\hat J) = \sqrt{J(1-J)/s}$}\\[4pt]
  plus de cases, moins d'erreur — en $1/\sqrt{s}$.
\end{frame}

\end{document}
